
\documentclass{article} % For LaTeX2e
\usepackage{iclr2025_conference,times}

% Optional math commands from https://github.com/goodfeli/dlbook_notation.
\input{math_commands.tex}

\usepackage{hyperref}
\usepackage{url}
\usepackage{booktabs}
\usepackage{tikz}


\title{Beyond Monolithic Memory: Regime-Aware Memory for Adaptive Agents}

% Authors must not appear in the submitted version. They should be hidden
% as long as the \iclrfinalcopy macro remains commented out below.
% Non-anonymous submissions will be rejected without review.

\author{Antiquus S.~Hippocampus, Natalia Cerebro \& Amelie P. Amygdale \thanks{ Use footnote for providing further information
about author (webpage, alternative address)---\emph{not} for acknowledging
funding agencies.  Funding acknowledgements go at the end of the paper.} \\
Department of Computer Science\\
Cranberry-Lemon University\\
Pittsburgh, PA 15213, USA \\
\texttt{\{hippo,brain,jen\}@cs.cranberry-lemon.edu} \\
\And
Ji Q. Ren \& Yevgeny LeNet \\
Department of Computational Neuroscience \\
University of the Witwatersrand \\
Joburg, South Africa \\
\texttt{\{robot,net\}@wits.ac.za} \\
\AND
Coauthor \\
Affiliation \\
Address \\
\texttt{email}
}

% The \author macro works with any number of authors. There are two commands
% used to separate the names and addresses of multiple authors: \And and \AND.
%
% Using \And between authors leaves it to \LaTeX{} to determine where to break
% the lines. Using \AND forces a linebreak at that point. So, if \LaTeX{}
% puts 3 of 4 authors names on the first line, and the last on the second
% line, try using \AND instead of \And before the third author name.

\newcommand{\fix}{\marginpar{FIX}}
\newcommand{\new}{\marginpar{NEW}}

%\iclrfinalcopy % Uncomment for camera-ready version, but NOT for submission.
\begin{document}


\maketitle

\begin{abstract}
Long-term memory lets agents reuse past experience, but semantic relevance does not imply validity under a new policy, goal, norm, or safety constraint.
We formalize these validity conditions as \emph{regimes} and study \emph{regime collapse}: cross-regime interference, memory homogenization, and minority-regime erasure caused by monolithic memory.
We propose \emph{RAM-minimal}, a concrete regime-aware memory mechanism with write-time regime assignment, conservative invariant promotion, and eligibility-first activation with hard incompatible-regime suppression.
We evaluate it on RegimeBench v0, a deterministic pilot benchmark with 720 queries across content moderation, robot planning, and cultural norm adaptation.
On the primary moderation domain, RAM-minimal reaches 96.0 RCA and 0.0 IMA, while a diagnostic metadata-filtering comparator reaches 94.3 RCA and 0.0 IMA; RAM-minimal further improves imbalanced-regime query preservation from 66.7 to 100.0.
Together, these results identify the limited settings where lifecycle validity changes retrieval outcomes beyond relevance ranking and read-time metadata filtering.
\end{abstract}
\section{Introduction}

Large language models are increasingly used as policy-conditioned moderation
systems. Instead of training a separate classifier for every platform rule,
jurisdiction, or customer requirement, a moderator can be prompted with a
natural-language policy and asked to judge whether a piece of content should be
allowed, contextualized, restricted, removed, or escalated. This interface is
attractive because moderation policies change frequently and because many
organizations need domain-specific rules that cannot be fully covered by a
static safety taxonomy.

However, deployment does not require only a policy-consistent judgment. It
requires an operational action. A policy may state that election misinformation
should be removed, that educational correction should be allowed, and that
borderline civic-integrity content may need context. These statements can be
reasonable at the level of institutional intent, while still underspecifying the
boundary that an operational moderator must apply. If a post quotes a false
ballot-validity rumor in order to debunk it, should the system allow it,
contextualize it, or remove it? If a satirical civic notice imitates an official
government format, is the decisive factor parody intent, deception risk, or the
requested user action? Such cases are not merely hard examples; they expose
whether the policy text contains enough information to determine the action
space it asks the model to execute.

We call this mismatch the \emph{policy-to-operation specification gap}: the gap
between a high-level moderation policy and the concrete operational decisions
that a model must make under that policy. The gap has two separable layers.
\emph{Boundary sufficiency} asks whether the policy specifies the variables that
separate decision families, such as correction versus endorsement, satire versus
impersonation, or prediction versus false certification. \emph{Action-granularity
sufficiency} asks whether the policy specifies the exact action within a
decision family, such as allow, contextualize, restrict, remove, or escalate.
The first layer concerns what kind of case the system is seeing; the second
concerns what intervention that case warrants.

This distinction matters because a model can appear to follow a policy while
remaining operationally unstable. A high-level policy may be sufficient for
coarse accept/reject judgments but insufficient for exact moderation actions.
Conversely, adding more examples or retrieving more precedents can create
pseudo-consensus around an underspecified boundary: the model becomes consistent
with the supplied cases without revealing whether those cases reflect the
current policy's intended operational boundary. Before building larger
moderation datasets, fine-tuning guardrail models, or adding long-term memory,
we therefore need a way to test whether the policy itself is sufficiently
specified for the decisions it delegates to the model.

We propose \emph{policy specification testing} as this missing diagnostic step.
Rather than generating synthetic training data, our method constructs matched
boundary probes that stress a policy's underspecified regions. Each probe is
paired with an expected decision family and an exact operational action. We then
compare model behavior under three policy views: an abstract policy, a
boundary-clarified policy that names the relevant distinctions, and a
boundary-plus-action policy that additionally defines the action rubric. The
goal is not to replace policy owners with a model-generated rulebook. The goal
is to reveal which boundary variables and action mappings are missing, so that
clarification can be targeted where it changes decisions.

A pilot study supports this decomposition. We build 12 action-granularity probes
for election moderation and evaluate two open models under the three policy
views. Moving from the abstract policy to the boundary-clarified policy reduces
decision-family error from 25.0\% to 8.3\% and instability from 58.3\% to
16.7\%. Adding the action rubric further reduces exact-action error from 54.2\%
to 33.3\% and action-granularity error from 45.8\% to 29.2\%. These numbers are
not intended as a benchmark-scale claim; they demonstrate the mechanism that
boundary clarification and action clarification repair different failure modes.

One representative case makes the mechanism concrete. For a post that quotes a
false ballot claim in order to debunk it and link official guidance, the
expected action is to contextualize rather than remove. Under the abstract
policy, both models remove the post. Under the boundary-clarified policy, both
models move to the correct decision family but choose allow. Under the
boundary-plus-action policy, both models choose contextualize. The same content
therefore moves through three distinct error regimes: wrong boundary, right
boundary but wrong action, and correct operational action.

This paper makes the following contributions:
\begin{itemize}
    \item We define the policy-to-operation specification gap in LLM-based
    content moderation: the failure mode where natural-language policy does not
    sufficiently determine operational moderation actions.
    \item We decompose the gap into boundary sufficiency and action-granularity
    sufficiency, separating coarse decision-family errors from exact-action
    errors.
    \item We introduce a lightweight policy specification testing protocol based
    on matched boundary probes, policy-view comparisons, instability analysis,
    and minimal clarification rescue.
    \item We provide pilot evidence that boundary clarification and action
    rubrics repair different classes of moderation failure, motivating policy
    specification testing as a preflight step before policy-conditioned judging,
    synthetic data generation, or memory-augmented moderation.
\end{itemize}

\section{Regime-Dependent Memory in Long-Term Agents}
\label{sec:problem_formulation}

Let $\mathcal{R}$ denote the set of regimes encountered by a long-term agent.
At time $t$, an agent receives input $x_t$, context $c_t$, active regime $r_t$, and memory state $M_t$.
A regime is not an arbitrary context tag: it is a condition that changes whether at least one memory may be reused.
We define the validity set of memory $m$ as
\begin{equation}
    V(m) \subseteq \mathcal{R}.
\end{equation}
Some memories are invariant across regimes; others are valid only under narrow policies, objectives, or norms.
This separates semantic relevance from regime validity:
\begin{equation}
    s_{rel}(m, x_t, c_t)
    \quad \text{and} \quad
    s_{val}(m, r_t).
\end{equation}
Optimizing $s_{rel}$ alone can activate memories that are topically close but invalid.

We use \emph{regime collapse} for failures caused by losing validity information during consolidation or retrieval.
The three measurable failures are \textbf{cross-regime interference}, where a memory from $r'\neq r$ steers behavior under $r$; \textbf{memory homogenization}, where distinct regime rules collapse into a generic heuristic; and \textbf{minority regime erasure}, where low-frequency regimes are overwritten by dominant regimes.
An effective memory system should therefore preserve regime-conditioned correctness, isolate incompatible regimes, share genuinely invariant memories, and retain rare regimes.

This formulation makes the evaluation target different from ordinary retrieval accuracy.
A retrieval system can select a topically relevant memory and still fail because the selected memory is valid only under a different regime.
Conversely, a system can avoid cross-regime interference by refusing to retrieve useful memories, which would reduce final task accuracy.
We therefore evaluate both final behavior and memory activation.
Regime-conditioned accuracy (RCA) measures whether the final answer matches the target under the active regime.
Cross-regime interference (CRI) measures whether the answer matches an explicitly invalid cross-regime negative.
Homogenization rate (HR) captures generic answers that avoid regime-specific commitments.
Minority-regime query preservation (MRP-Acc) measures downstream accuracy on rare-regime queries, and invalid memory activation (IMA) measures the fraction of retrieved memories whose assignment is incompatible with the active regime.
These metrics are deliberately redundant: a credible method should reduce CRI and IMA without simply becoming generic or suppressing all useful memory.

The active regime can be explicit, inferred, or partially specified.
This paper studies the explicit-regime case because it isolates memory lifecycle validity from regime inference.
That choice makes the closest comparator unusually strong: if a method already receives active-regime metadata, a simple read-time metadata filter can prevent many invalid activations.
The remaining question is narrower but still important for long-term agents: whether write-time store separation, conservative invariant promotion, and retention pressure change outcomes beyond read-time filtering alone.

Formally, a memory method receives a construction stream $\mathcal{H}=\{e_1,\ldots,e_T\}$ and produces a memory state $M_T$.
At query time it receives $(x,c,r)$ and returns a prediction $\hat y$ after retrieving a context $C_K(M_T,x,c,r)$ of at most $K$ memories.
The usual relevance objective encourages $C_K$ to contain memories with high semantic similarity to $(x,c)$.
Regime-dependent memory adds a constraint: activated memories should be compatible with $r$, and memories that are useful across regimes should be shared only when there is evidence for that invariance.
This turns memory construction into part of the algorithm rather than a passive log of previous examples.

The failure taxonomy also clarifies what would falsify the proposed method.
If semantic retrieval has low IMA and low CRI on conflict pairs, then the relevance-validity gap is not present in the benchmark.
If metadata filtering matches RAM-minimal on minority-regime retention and invariant transfer, then write-time lifecycle structure is unnecessary under the tested conditions.
If RAM-minimal reduces IMA only by failing to retrieve useful memories, RCA should drop.
These failure conditions are important because the paper's claim is intentionally narrow: RAM-minimal should improve exactly the pressure points where lifecycle validity matters, not dominate every explicit-regime memory comparator.

\section{RAM-minimal}
\label{sec:methodology}

RAM-minimal is the concrete method studied in this paper.
It instantiates the RAM principle with the smallest set of operations needed to test the relevance-validity gap: write-time regime assignment, conservative invariant promotion, and eligibility-first activation with incompatible-regime suppression.
The method assumes that RegimeBench v0 provides an explicit active regime for each construction and query example.
Latent regime inference, learned controllers, deprecation, split, and merge are useful extensions, but they are not required for the main empirical claim.

Each memory record is
\begin{equation}
    m_i = (u_i, z_i, a_i, v_i, \eta_i, \sigma_i),
\end{equation}
where $u_i$ is a stable memory identifier, $z_i$ is the memory content, $a_i \in \{0, ?\}\cup\mathcal{R}$ is its assignment, $v_i$ stores the current validity scope, $\eta_i$ stores provenance such as source example, timestamp, feedback, and rationale, and $\sigma_i$ is the status.
RAM-minimal decomposes memory into invariant, regime-specific, and optional uncertain stores:
\begin{equation}
    M_t = M_t^0 \cup \bigcup_{r \in \mathcal{R}} M_t^r \cup M_t^?,
\end{equation}
where $M_t^0$ contains invariant memories and $M_t^r$ contains memories valid under regime $r$.
The store decomposition is the main design choice.
A flat memory pool can attach labels to memories, but it still consolidates all regimes through the same capacity pressure and update path.
RAM-minimal instead lets each regime maintain local evidence while reserving $M^0$ for memories that have explicit cross-regime support.
This distinction matters most when one regime is rare: a low-frequency memory can be overwritten in a flat pool even if read-time filtering is later perfect.

Given a construction example with explicit regime $r_t$, RAM-minimal writes the extracted memory to $M^{r_t}$ by default.
Write-time assignment is conservative: a memory observed under one regime is not treated as globally valid.
It promotes a rule group $g$ to $M^0$ only when positive evidence appears in at least $k$ distinct regimes; v0 uses $k=2$ and deterministic rule identifiers.
Promotion is therefore a high-precision operation rather than an aggressive sharing rule.
In a richer implementation, the rule group could be replaced by a learned equivalence class over semantically similar memories, but the reported v0 avoids that additional modeling assumption.

At query time, RAM-minimal first applies hard eligibility:
\begin{equation}
    \mathcal{E}(r)=M^0 \cup M^r.
\end{equation}
Eligible memories are then ranked by relevance plus an optional validity bonus:
\begin{equation}
    \mathrm{Score}(m,x,c,r)=s_{rel}(m,x,c)+\gamma s_{val}(m,r).
\end{equation}
The reported implementation uses TF-IDF cosine similarity, $K=3$, $\gamma=1.0$, and hard suppression of memories assigned to incompatible regimes; ablations show that hard eligibility, not the soft validity bonus, drives the main v0 effect.
Full pseudocode and scoring details are in Appendix~\ref{app:implementation_details}.

\begin{table}[t]
\centering
\footnotesize
\begin{tabular}{p{0.95\linewidth}}
\toprule
\textbf{Algorithm 1: RAM-minimal, compact form} \\
\midrule
Initialize invariant store $M^0$ and regime stores $\{M^r:r\in\mathcal{R}\}$. \\
For each construction example $(x,c,r,y,\mathrm{rule})$, write the extracted memory to $M^r$. \\
Promote a rule group to $M^0$ only if it has positive evidence in at least $k$ regimes. \\
For query $(x,c,r)$, form eligible memories $\mathcal{E}(r)=M^0\cup M^r$. \\
Rank eligible memories by $s_{rel}(m,x,c)+\gamma s_{val}(m,r)$ and retrieve top $K$. \\
Predict using the retrieved memories and the same deterministic task policy as all comparators. \\
\bottomrule
\end{tabular}
\caption{Main-text summary of the proposed method. All methods receive the same construction/query streams and $K=3$ retrieval budget.}
\label{alg:ram_compact}
\end{table}

Regime-Tagged Retrieval serves as the diagnostic metadata-filtering comparator because it also receives active regime metadata.
However, it treats regime labels as read-time filters over a flat memory pool.
RAM-minimal changes the memory lifecycle before retrieval: write location is regime-specific, invariant memories are promoted under conservative evidence, incompatible regimes are made ineligible under explicit active regimes, and minority-regime retention is an explicit diagnostic target.
The distinction is empirical, not merely terminological.
When Regime-Tagged Retrieval matches RAM-minimal on minority and invariant-sharing diagnostics, lifecycle construction adds little in that setting.

\begin{table*}[t]
\centering
\footnotesize
\begin{tabular}{p{0.24\linewidth} p{0.31\linewidth} p{0.33\linewidth}}
\toprule
\textbf{Dimension} & \textbf{Regime-Tagged Retrieval} & \textbf{RAM-minimal} \\
\midrule
Write-time separation & Flat pool with labels & Separate invariant and regime-specific stores \\
\midrule
Invariant memory & Not explicit & Dedicated $M^0$ store \\
\midrule
Promotion & No lifecycle promotion & Conservative promotion with cross-regime support \\
\midrule
Incompatible regime suppression & Read-time filter & Eligibility control before ranking plus optional validity bonus \\
\midrule
Minority-regime preservation & Not a design target & Explicit retention and query diagnostic \\
\midrule
Evolving validity & Weak static labels & Extension point through status and provenance \\
\bottomrule
\end{tabular}
\caption{RAM-minimal differs from regime-tagged retrieval by changing memory construction and eligibility, not only read-time label filtering.}
\label{tab:ram_vs_tagged}
\end{table*}

The proposed method is intentionally minimal.
It does not learn a controller, infer latent regimes, or implement deprecation in the reported experiments.
Those operations belong to the broader RAM family, but including them in the main claim would make the paper harder to falsify.
The empirical test here is whether the smallest lifecycle intervention--separate regime stores, conservative invariant promotion, and hard eligibility--is enough to expose failures that relevance-only retrieval and flat metadata filtering hide.

\paragraph{Write rule.}
The write rule is deliberately conservative.
For an experience $e_t=(x_t,c_t,r_t,y_t,\mathrm{feedback}_t)$, RAM-minimal extracts a memory containing the observed behavior and writes it to $M^{r_t}$ unless it already belongs to a promoted invariant group.
This prevents a single observation from becoming a global rule.
The v0 implementation uses deterministic rule identifiers supplied by the benchmark, so the experiment does not claim to solve memory clustering.
This is a limitation, but it also makes the mechanism easy to audit: when a memory is promoted, the artifact records the source rule and the regimes that support it.

\paragraph{Promotion rule.}
Promotion to $M^0$ requires evidence in at least $k$ distinct regimes and an invariant-candidate flag in the construction data.
The default is $k=2$.
This rule is intentionally asymmetric: it is easier to keep a memory regime-specific than to promote it globally.
The asymmetry matters because false promotion is costly in a regime-aware system: a memory incorrectly placed in $M^0$ becomes eligible everywhere.
In the reported pilot, promotion has a small but positive diagnostic effect; it is not the main driver of performance.

\paragraph{Activation rule.}
At query time, RAM-minimal first forms an eligible set $\mathcal{E}(r)=M^0\cup M^r$ and only then ranks memories.
This is the central difference from relevance-only retrieval, where all memories compete in the same ranking space.
It is also different from a flat tag filter because the memory was already written under regime-specific capacity pressure.
The optional validity bonus $\gamma s_{val}$ is included for completeness, but the primary moderation ablation shows that setting $\gamma=0$ leaves the main results unchanged.
We therefore interpret the v0 evidence as support for eligibility and store separation, not for a learned validity scorer.

\paragraph{Capacity diagnostic.}
Because separate stores can retain more memories than a flat store, capacity could explain the minority-regime result.
To test this, we add a tagged-capacity sensitivity diagnostic.
RAM-minimal is first built normally and then pruned by removing the oldest memories from the largest stores until its realized total memory count matches Regime-Tagged Retrieval for the same domain and split.
This diagnostic is not a deployable variant; it is a stress test for the artifact.
In v0, the tagged-capacity version preserves the moderation and imbalanced-regime gains, suggesting that the observed minority-retention effect is not only a consequence of storing more memories.

\section{RegimeBench v0}
\label{sec:regimebench_v0}

RegimeBench v0 evaluates long-term agent tasks where similar inputs require different behavior under different regimes.
The central question is whether a memory system can reuse experience without applying it outside its validity conditions.
Each query instance is
\begin{equation}
    (x_i, c_i, r_i, y_i, \mathcal{Y}_i^{-}),
\end{equation}
where $x_i$ is the query or environment state, $c_i$ is metadata, $r_i$ is the active regime, $y_i$ is the regime-correct output, and $\mathcal{Y}_i^{-}$ contains plausible outputs valid under other regimes but invalid under $r_i$.
Each example also stores a rule identifier, rationale, invariant-candidate flag, and minority-regime flag.
The benchmark has four splits: seen-regime recall, imbalanced-regime retention, conflict-pair interference, and evolving-regime diagnostics.
Content moderation is the primary domain; robot planning and cultural norm adaptation test transfer to procedural and social regimes.

RegimeBench v0 is deliberately small enough to inspect and rerun.
It contains 600 balanced construction prototypes and 720 query examples.
Because each split is evaluated with its own construction stream, the released \texttt{dataset.jsonl} logs 2,085 split-specific construction events plus the 720 queries.
The released \texttt{data\_summary.json} records 2,805 total rows.
All reported numbers are computed from \texttt{predictions.jsonl}, \texttt{retrieved\_memories.jsonl}, \texttt{memory\_states.jsonl}, \texttt{metrics.json}, \texttt{diagnostics.json}, \texttt{ablations.json}, \texttt{capacity\_matched.json}, and \texttt{memory\_costs.json}.

\begin{table*}[t]
\centering
\scriptsize
\begin{tabular}{@{}lcccc@{}}
\toprule
\textbf{Domain} & \textbf{Regimes} & \textbf{Balanced prototypes} & \textbf{Queries} & \textbf{Conflict queries} \\
\midrule
Content moderation & 5 & 250 & 300 & 100 \\
Robot planning & 3 & 150 & 180 & 60 \\
Cultural norms & 4 & 200 & 240 & 80 \\
\midrule
Total & 12 & 600 & 720 & 240 \\
\bottomrule
\end{tabular}
\caption{RegimeBench v0 scale. The artifact logs 2,085 split-specific construction events because seen, imbalanced, conflict, and evolving splits use separate construction streams.}
\label{tab:regimebench_scale}
\end{table*}

Examples are constructed from deterministic templates.
Each template specifies a base situation, surface variants, active regime, target output, cross-regime negative outputs, rule identifier, rationale, and invariant-candidate flag.
Construction and query splits do not share exact wording.
Conflict pairs preserve semantic similarity while changing the active regime and target output.
Active-regime metadata is made available fairly: methods that use explicit regimes receive the same active regime at query time, while relevance-only methods are reported as not using it.
Appendix~\ref{app:regimebench_construction} gives construction, leakage-control, and judge-protocol details.

The four splits are designed to separate different failure modes.
The seen-regime split checks whether methods can reuse memories when construction and query regimes match.
The conflict-pair split creates semantically similar inputs whose correct outputs differ across regimes, making relevance-only retrieval brittle by construction.
The imbalanced-regime split down-samples a designated minority regime during construction and oversamples it at query time, exposing whether memory capacity and consolidation preserve rare regimes.
The evolving-regime split changes selected labels between construction and query.
Because RAM-minimal does not implement deprecation or split/merge, evolving-regime results are diagnostic rather than central evidence for the method.

Evaluation is exact-match in v0.
For moderation and robot planning, outputs are compact labels such as \textit{allow}, \textit{restrict}, or a discrete robot action.
For cultural norm adaptation, deterministic response templates play the same role.
This sacrifices natural-language realism in exchange for full traceability: every error can be attributed to a target output, a cross-regime negative set, and the retrieved memories.
A future free-form benchmark would need fixed judge prompts and judge calibration, but those judge outputs are not used for any reported number in this paper.

\paragraph{Domain design.}
The moderation domain is the primary testbed because policy regimes make conditional validity easy to inspect.
The same surface post can be allowed under a public-discussion regime, restricted under a protected-group-harm regime, or handled differently during a crisis event.
Robot planning turns the same idea into procedural objectives: a route or action can be appropriate under efficiency-first behavior but invalid when safety-first constraints are active.
Cultural norm adaptation tests social validity: a direct response may be appropriate under an informal-peer regime but invalid under a hierarchy-sensitive regime.
These domains are not meant to exhaustively cover real policies, robotics, or cultural behavior.
They are controlled probes for the same abstract memory failure.

\paragraph{Leakage controls.}
Construction and query examples share rule identities but not exact wording.
Conflict-pair examples preserve semantic similarity while changing the active regime, target output, and negative set.
This prevents methods from succeeding by memorizing exact strings while still making relevance-only retrieval vulnerable to topically close invalid memories.
All methods receive the same construction stream in the same order.
Methods that use explicit regimes receive the same active-regime metadata; methods that do not use it are reported separately so that the comparison does not hide the advantage of regime information.

\paragraph{Pilot status.}
RegimeBench v0 is a deterministic controlled pilot; a full-scale benchmark remains future work.
Its value is that it makes memory-lifecycle failures easy to reproduce and inspect.
Its weakness is that deterministic labels and TF-IDF retrieval are cleaner than real long-term agent memory.
The experiments therefore support a mechanism claim, not a deployment claim: if an agent stores conditionally valid memories in a monolithic state, relevance alone can activate invalid memories, and lifecycle-aware storage can prevent a subset of those failures.

\section{Experiments}
\label{sec:experiments}

The experiments test whether the proposed method reduces regime collapse.
We compare against No Memory, Monolithic Memory, Semantic Retrieval, Context-Aware Retrieval, and Regime-Tagged Retrieval, a diagnostic metadata-filtering comparator.
All results are computed from deterministic RegimeBench v0 predictions, retrieved memories, and memory states saved under \texttt{results/regimebench\_v0/}.
All retrieval methods use the same top-$K=3$ activation budget; because lifecycle stores can retain different numbers of memories, we also report memory cost and a tagged-capacity sensitivity test.

We test four falsifiable claims.
H1: relevance-only retrieval increases invalid activation on conflict pairs.
H2: metadata filtering is strong when regimes are explicit but weaker under imbalanced minority pressure.
H3: RAM-minimal reduces CRI and IMA without reducing RCA.
H4: invariant promotion has a small positive diagnostic effect rather than a broad standalone result.
Metrics are RCA, CRI, HR, minority-regime query preservation (MRP-Acc), IMA, and cost.
IMA is only defined for methods that retrieve individual assigned memories; monolithic summaries are evaluated by CRI and HR instead.
Here, cost in \texttt{metrics.json} means activation cost, the average number of retrieved memories per query; stored-memory cost is reported separately in \texttt{memory\_costs.json}.

RCA, CRI, and HR are computed from the final prediction.
RCA is exact-match accuracy against the active-regime target.
CRI checks whether the prediction falls in the cross-regime negative set $\mathcal{Y}^{-}$, and HR checks whether the prediction collapses to a generic domain-specific response.
MRP-Acc is computed only on minority-regime queries, so it measures downstream preservation rather than just memory count.
IMA is computed from retrieved memory assignments: a retrieved memory is invalid if it is neither invariant nor assigned to the active regime.
This metric is not defined for Monolithic Memory because its summaries do not expose individual regime assignments; we therefore mark Monolithic IMA as N/A in the main table.
Table~\ref{tab:main_results_template} asks whether invalid activation can be reduced without sacrificing regime-conditioned accuracy.

\begin{table*}[t]
\centering
\small
\begin{tabular}{lccccc}
\toprule
\textbf{Method} &
\textbf{RCA $\uparrow$} &
\textbf{CRI $\downarrow$} &
\textbf{HR $\downarrow$} &
\textbf{MRP-Acc $\uparrow$} &
\textbf{IMA $\downarrow$} \\
\midrule
No Memory & 16.7 & -- & 83.3 & -- & -- \\
Monolithic Memory & 34.3 & 65.7 & 64.7 & 23.2 & N/A \\
Semantic Retrieval & 44.0 & 56.0 & 0.0 & 14.5 & 81.8 \\
Context-Aware Retrieval & 51.0 & 49.0 & 0.0 & 36.2 & 69.4 \\
Regime-Tagged Retrieval & 94.3 & 5.7 & 0.0 & 92.8 & 0.0 \\
RAM-minimal & 96.0 & 4.0 & 0.0 & 100.0 & 0.0 \\
\bottomrule
\end{tabular}
\caption{Primary content moderation results on RegimeBench v0. MRP-Acc is accuracy on minority-regime queries. Monolithic IMA is N/A because monolithic summaries do not retain individual regime assignments; their failures are captured by CRI and HR.}
\label{tab:main_results_template}
\end{table*}

Bootstrap intervals over moderation queries support the same pattern.
RAM-minimal obtains RCA $96.0$ with a 95\% bootstrap interval of $[93.7,98.0]$, compared with $94.3$ $[91.7,97.0]$ for Regime-Tagged Retrieval and $44.0$ $[38.7,49.7]$ for Semantic Retrieval.
Against relevance-only retrieval, the strongest diagnostic gap is IMA: Semantic Retrieval activates invalid memories in $81.8$ $[77.6,85.9]$ percent of retrieved contexts, while RAM-minimal and Regime-Tagged Retrieval both remain at $0.0$.
The main table shows two points: read-time metadata filtering is strong under explicit regimes, and RAM-minimal adds minority retention and invariant-store benefits on top of that diagnostic comparator.

Table~\ref{tab:split_results} reports split-level diagnostics aligned with the hypotheses.
It asks where explicit-regime metadata filtering is sufficient and where lifecycle construction changes the outcome.
Regime-Tagged Retrieval is strong on seen regimes and eliminates conflict-pair interference, but RAM-minimal improves minority-regime query preservation through separate regime stores and conservative invariant promotion.
The evolving-regime split is intentionally diagnostic: RAM-minimal and Regime-Tagged Retrieval tie in this deterministic v0, so the present evidence for RAM-minimal over metadata filtering comes from imbalanced-regime retention and small invariant/promotion diagnostics rather than from evolving-regime gains.
Across transfer domains, RAM-minimal obtains 92.8 RCA on robot planning and 93.8 on cultural norm adaptation, compared with 87.2 and 89.6 for Regime-Tagged Retrieval.

Memory capacity is a possible confound, since the default RAM-minimal setting retains invariant and regime-specific stores.
Table~\ref{tab:capacity_diagnostic} therefore reports stored-memory cost and a tagged-capacity sensitivity in which RAM-minimal is pruned to the same realized total memory count as Regime-Tagged Retrieval for each domain and split.
The moderation and imbalanced-regime pattern is unchanged in this controlled diagnostic, suggesting that the v0 minority-retention gain is not explained only by storing more memories.

\begin{table*}[t]
\centering
\footnotesize
\begin{tabular}{lcccc}
\toprule
\textbf{Method} &
\textbf{Seen RCA $\uparrow$} &
\textbf{Conflict CRI $\downarrow$} &
\textbf{Imbal. MRP-Acc $\uparrow$} &
\textbf{Evolving RCA $\uparrow$} \\
\midrule
Semantic Retrieval & 59.4 & 66.7 & 33.3 & 15.6 \\
Context-Aware Retrieval & 64.4 & 58.8 & 42.9 & 21.2 \\
Regime-Tagged Retrieval & 99.4 & 0.0 & 66.7 & 75.0 \\
RAM-minimal & 100.0 & 0.0 & 100.0 & 75.0 \\
\bottomrule
\end{tabular}
\caption{Split-level diagnostics on RegimeBench v0. The strongest RAM-minimal advantage over metadata filtering is imbalanced minority-regime query preservation; evolving-regime behavior remains tied in this v0 pilot.}
\label{tab:split_results}
\end{table*}

\begin{table*}[t]
\centering
\scriptsize
\setlength{\tabcolsep}{4pt}
\begin{tabular}{@{}lccccc@{}}
\toprule
\textbf{Method} &
\textbf{Mean mem.} &
\textbf{Max mem.} &
\textbf{Inv. mem.} &
\textbf{Mod. RCA} &
\textbf{Imbal. MRP-Acc} \\
\midrule
Regime-Tagged & 95.0 & 125 & 0.0 & 94.3 & 66.7 \\
RAM-minimal & 158.8 & 225 & 34.5 & 96.0 & 100.0 \\
RAM-minimal (tagged cap.) & 95.0 & 125 & 18.7 & 96.0 & 100.0 \\
\bottomrule
\end{tabular}
\caption{Memory-cost and capacity sensitivity from \texttt{memory\_costs.json} and \texttt{capacity\_matched.json}. The tagged-capacity variant is a diagnostic version of the proposed method pruned to the same realized total memory count as Regime-Tagged Retrieval.}
\label{tab:capacity_diagnostic}
\end{table*}

The diagnostic degradation localizes most components: removing write-time assignment collapses to relevance-only behavior, removing suppression increases invalid activation, and removing invariant memory or promotion slightly weakens transfer.
Setting $\gamma=0$ has no measurable effect in this explicit-regime v0, indicating that hard eligibility control rather than soft validity weighting drives the main gain.
The diagnostic contrasts from \texttt{diagnostics.json} summarize this narrow pattern: RAM-minimal is +1.7 RCA over Regime-Tagged Retrieval on moderation, +33.3 MRP-Acc on imbalanced queries, and tied on evolving-regime RCA; Semantic Retrieval has +66.7 conflict CRI relative to RAM-minimal.
Table~\ref{tab:ablation_results_template} asks which operations account for this pattern.

\begin{table}[t]
\centering
\small
\begin{tabular}{lccc}
\toprule
\textbf{Variant} &
\textbf{RCA $\uparrow$} &
\textbf{CRI $\downarrow$} &
\textbf{IMA $\downarrow$} \\
\midrule
RAM-minimal & 96.0 & 4.0 & 0.0 \\
w/o Write-Time Assignment & 44.0 & 56.0 & 79.9 \\
w/o Regime Suppression & 91.0 & 9.0 & 3.0 \\
w/o Invariant Memory & 94.3 & 5.7 & 0.0 \\
w/o Promotion & 94.3 & 5.7 & 0.0 \\
$\gamma=0$ & 96.0 & 4.0 & 0.0 \\
\bottomrule
\end{tabular}
\caption{Primary moderation ablations on RegimeBench v0. The strongest evidence supports write-time assignment and suppression; invariant memory and promotion have a small positive diagnostic effect, while the soft validity bonus has no measurable effect in this explicit-regime pilot.}
\label{tab:ablation_results_template}
\end{table}

The ablations sharpen the interpretation of H4.
Removing write-time assignment reduces moderation RCA from 96.0 to 44.0 and raises IMA to 79.9, closely resembling relevance-only retrieval.
Removing suppression preserves much of the method but increases CRI and IMA, showing that store separation alone is insufficient if incompatible stores remain eligible.
Removing invariant memory or promotion drops RCA from 96.0 to 94.3, matching Regime-Tagged Retrieval; this is a positive but small diagnostic effect, not evidence for a broad learned promotion mechanism.
Finally, $\gamma=0$ leaves the main numbers unchanged, so the current evidence supports hard eligibility and lifecycle storage more strongly than soft validity weighting.

Failure traces support the aggregate pattern.
In conflict-pair queries, Semantic Retrieval has CRI $66.7$ and IMA $74.6$, while RAM-minimal has CRI $0.0$ and IMA $0.0$.
For a public-discussion query about flooding an event page with angry comments, Semantic Retrieval predicts \textit{restrict} after retrieving an incompatible protected-group-harm memory; Regime-Tagged Retrieval and RAM-minimal both predict \textit{allow}.
This case is a seen-regime relevance-validity failure rather than the source of the conflict-pair aggregate.

The transfer domains show the same qualitative pattern but should not be overread.
On robot planning and cultural norm adaptation, RAM-minimal reaches 92.8 and 93.8 RCA, respectively, while Regime-Tagged Retrieval reaches 87.2 and 89.6.
These domains are useful because they show that regime-dependent memory is not only a moderation artifact; however, they are still deterministic template domains.
The transfer results are best read as additional controlled evidence: relevance-only retrieval can activate invalid memories, explicit-regime metadata filtering is a strong comparator, and RAM-minimal adds a narrow lifecycle benefit concentrated on minority retention and conservative invariant sharing.

\section{Discussion}
\label{sec:discussion}

The experiments support a narrow mechanism claim rather than a general replacement for retrieval, metadata filters, or learned memory controllers.
When memories are conditionally valid, relevance-only retrieval can activate invalid memories; in the controlled v0 setting, a minimal lifecycle mechanism prevents those activations without reducing task accuracy.
Regime-Tagged Retrieval is an intentionally strong explicit-regime comparator: it receives the active regime and should remove most invalid activations when labels are reliable.
RAM-minimal matters only where writing, promotion, and capacity pressure change what survives until query time.
The capacity and ablation diagnostics make this interpretation more precise.
The tagged-capacity sensitivity keeps realized mean stored memory at 95.0 for both Regime-Tagged Retrieval and RAM-minimal (tagged capacity), while RAM-minimal still obtains 100.0 imbalanced MRP-Acc, so the v0 minority gain is not explained only by larger storage.
The ablations identify write-time assignment and hard suppression as the main components; invariant memory and promotion provide a small positive diagnostic effect, and the soft validity bonus has no measurable effect under explicit regimes.
The current evidence therefore supports eligibility before relevance ranking, not a learned validity scorer or full lifecycle controller.

RegimeBench v0 intentionally favors inspection over realism.
The template-based domains make conflict pairs, negative sets, and retrieved-memory traces easy to audit.
This is useful for debugging memory mechanisms, but it also means that natural-language ambiguity, latent regime inference, and judge calibration remain open.
A stronger future benchmark should add free-form generation and calibrated judges, but the role of v0 is to isolate a failure that larger evaluations can otherwise hide.
The practical implication is modest: when regimes are static, explicit, and balanced, read-time metadata filtering may be sufficient; RAM-minimal is most useful when regimes are imbalanced, invariant memories should transfer, or update rules must preserve validity boundaries over time.

\section{Related Work}
\label{sec:related_work}

Long-term memory systems for language agents store, retrieve, summarize, and reason over experiences across sessions~\citep{park2023generative,packer2023memgpt,sumers2023coala,wu2024longmemeval}.
These works primarily ask whether agents can remember and use past information.
RAM-minimal targets a complementary failure: the agent remembers a relevant memory, but that memory is no longer valid under the active policy, objective, norm, or safety constraint.
The distinction is not that prior memory systems ignore retrieval quality; rather, their memory states usually do not make conditional validity a first-class lifecycle property.
If all memories are consolidated into a single pool, the system can improve recall while still losing the boundary conditions under which a memory should be reused.

Context-aware and intent-aware retrieval improve read-time selection by conditioning on task context, interaction state, or downstream utility~\citep{xu2025amem,utilityrag2025,influencecontext2025}.
The closest diagnostic alternative is simple metadata filtering: store memories with observed regime labels and filter the flat pool by the active regime.
We instantiate this strong comparator as Regime-Tagged Retrieval in our experiments.
The proposed method differs only where lifecycle validity matters: write-time store separation, conservative invariant promotion, and minority-regime retention.
When active regimes are explicit, metadata filtering should be strong, and the experiments confirm that it eliminates most invalid activation.
The remaining question is whether read-time labels are enough under memory construction pressure, rare regimes, and invariant sharing.

Memory update, compression, and forgetting methods decide what an agent should retain, summarize, overwrite, or discard as interactions accumulate.
These mechanisms are complementary to RAM-minimal: they can reduce stale or redundant memories, but they do not by themselves specify whether a retained memory remains valid under a different policy, norm, objective, or safety constraint.
Regime-aware memory can therefore act as a lifecycle constraint on top of update or forgetting policies.
Update rules decide whether a memory should remain in the system; regime validity decides where it is stored and when it is eligible for activation.

Pluralistic alignment and cultural adaptation study systems that support diverse values, preferences, and norms~\citep{sorensen2024roadmap,culturalpalette2024,valuepluralism2025}.
RAM-minimal studies the memory-layer analogue: whether a single agent preserves distinct norm- or policy-specific memories during consolidation rather than averaging them into generic heuristics.
It is not a multi-agent or mixture-of-experts method; regimes are validity structures attached to memories, not separate experts.
This makes the proposed method narrower than pluralistic alignment: it does not decide which value system should be used, and it does not infer social context from scratch.
Instead, given an active regime, it asks whether long-term memory preserves the distinctions needed to act correctly under that regime.

\section{Limitations}
\label{sec:limitations}

The reported experiments are a deterministic RegimeBench v0 pilot, not a full deployment benchmark.
It uses templated examples, exact-match labels, TF-IDF retrieval, and deterministic task policies so that every result can be traced to predictions, retrieved memories, and memory states.
This design isolates the relevance-validity gap but overstates the cleanliness of real agent memory.
RAM-minimal also assumes explicit active regimes in the main evaluation and therefore does not solve latent regime inference.
When regimes are explicit, balanced, and static, metadata filtering can be nearly sufficient; RAM-minimal is most useful when invariant transfer, minority retention, or lifecycle validity boundaries matter.
Finally, the current implementation uses deterministic rule identifiers for promotion, making invariant discovery easier than learning equivalence classes from natural language.

\section{Conclusion}
\label{sec:conclusion}

We introduced RAM-minimal, a concrete regime-aware memory mechanism for long-term agents whose memories are conditionally valid under goals, norms, policies, or task constraints.
RegimeBench v0 shows that relevance-only retrieval can activate invalid memories, while RAM-minimal matches strong metadata filtering on invalid activation and improves minority-regime query preservation in the controlled pilot.
The broader lesson is simple: memory-augmented agents should not merely remember more; they should remember the regimes under which their memories remain valid.


\bibliography{iclr2025_conference}
\bibliographystyle{iclr2025_conference}

\appendix
\section{Implementation Details for RAM-minimal and Extensions}
\label{app:implementation_details}

\subsection{Validity Scoring and Normalization}
\label{app:validity_scoring}

RAM-minimal separates semantic relevance from regime validity.
In the explicit-regime v0 implementation, hard eligibility is applied before ranking; for eligible memories, the implementation can also compute a relevance score $s_{rel}(m,x,c)$ and a validity score $s_{val}(m,r)$.
When regimes are explicit, $s_{val}$ is computed from memory assignment:
\begin{equation}
    s_{val}(m,r)
    =
    \mathbb{I}[a(m)=0]
    +
    \mathbb{I}[a(m)=r],
\end{equation}
where $a(m)=0$ denotes invariant memory.
Invariant memories and memories assigned to the active regime receive high compatibility; memories assigned to other regimes are suppressed.

In future variants where regimes are latent, noisy, or partially specified, $s_{val}$ could be estimated by a compatibility model:
\begin{equation}
    s_{val}(m,r) = p_\psi(\mathrm{valid}=1 \mid m,r,x,c),
\end{equation}
where $p_\psi$ may be a classifier, an instruction-tuned language model, or an LLM judge prompted to decide whether memory $m$ remains valid under regime $r$.
For uncertain regimes, RAM marginalizes over the inferred regime distribution:
\begin{equation}
    s_{val}(m)
    =
    \sum_{r \in \mathcal{R}}
    q_\phi(r \mid x,c,h) s_{val}(m,r).
\end{equation}

In larger variants, relevance and validity scores can be normalized to $[0,1]$ on a validation split and combined after eligibility control.
The retrieval score is
\begin{equation}
    \mathrm{Score}(m,x,c,r)
    =
    \mathrm{Norm}(s_{rel}(m,x,c))
    +
    \gamma \mathrm{Norm}(s_{val}(m,r)).
\end{equation}
The coefficient $\gamma$ controls the optional validity bonus.
The main v0 setting fixes $\gamma=1.0$, but the $\gamma=0$ ablation is unchanged on the primary moderation domain; this is why the main claim emphasizes store separation and hard eligibility rather than learned validity scoring.
Tuning $\gamma$ on validation regimes is an extension for larger benchmarks.

\subsection{Memory Actions}
\label{app:memory_actions}

The broader RAM family supports several memory actions, but RAM-minimal requires only writing, activation, suppression, and promotion.
We describe practical decision rules below.

\paragraph{Writing.}
Given a new experience $e_t$, RAM extracts a candidate memory $m_t$ and assigns it to either a regime-specific memory $M^r$ or invariant memory $M^0$.
If the active regime is explicit, the default action is \textsc{WriteToRegime}$(r_t)$.
If the active regime is latent in a future extension, RAM can write to the most likely regime $\hat r_t=\arg\max_r q_\phi(r \mid x_t,c_t,h_t)$ when confidence exceeds a threshold; otherwise, the memory is marked uncertain and revisited after additional evidence.

\paragraph{Promotion to invariant memory.}
A memory is promoted from $M^r$ to $M^0$ only when there is evidence that it is valid across regimes.
Promotion occurs in three cases:
(1) the data construction or a heuristic explicitly marks the memory as regime-independent;
(2) the same memory or rule receives positive feedback under multiple regimes;
or (3) in a future latent-regime variant, a compatibility model predicts high $s_{val}(m,r)$ for several distinct regimes.
Promotion therefore avoids treating a single-regime success as globally valid.

\paragraph{Suppression.}
During retrieval, memories assigned to regimes incompatible with the inferred active regime are suppressed by validity scoring.
Suppression is soft when the regime distribution is uncertain and hard when the regime is explicit and high confidence.
RAM can therefore reduce cross-regime interference without discarding memories that may become useful later.

\paragraph{Deprecation.}
A memory is deprecated when it repeatedly causes incorrect behavior, receives negative feedback, or conflicts with an updated regime definition.
Deprecated memories are retained with provenance metadata but excluded from normal retrieval.
This is useful for evolving-regime settings, where policy updates may invalidate memories that were previously correct.

\paragraph{Split and merge.}
Split and merge are optional lifecycle actions.
\textsc{SplitRegime} is used when a regime contains internally inconsistent memories that lead to systematic errors.
\textsc{MergeRegime} is used when two regimes show high mutual compatibility and separating them produces unnecessary duplication.
These actions are not required for the minimal RAM implementation but are useful when regime taxonomies are initially coarse or change over time.

\paragraph{Oracle diagnostic.}
The code also logs an \emph{Oracle Regime Memory} diagnostic that receives ground-truth regime labels and exact rule identity when ranking memories.
It is not a deployable method and is omitted from the main tables because deterministic v0 makes RAM-minimal match it on several splits.
Its purpose is only to check whether remaining errors come from memory mechanics rather than task ambiguity.

\section{RegimeBench Construction Details}
\label{app:regimebench_construction}

RegimeBench is designed to evaluate whether memory systems preserve conditional validity across long-term memory construction and query phases.
All domains follow the same construction pattern.
The reported v0 pilot is deterministic: examples are generated from fixed templates, retrieval uses TF-IDF cosine similarity, and task outputs are exact-match labels.

\subsection{Memory Stream Construction}

For each domain, we define a set of regimes $\mathcal{R}$ and construct a memory stream
\begin{equation}
    \mathcal{H} = \{e_1,\ldots,e_T\}.
\end{equation}
Each experience contains an input, context metadata, active regime, target behavior, and feedback or rationale:
\begin{equation}
    e_t=(x_t,c_t,r_t,y_t,\mathrm{feedback}_t).
\end{equation}
The memory stream is processed sequentially, so memory systems cannot revise earlier decisions unless their lifecycle update rules allow them to do so.
All compared methods receive the same stream in the same order.
In the balanced construction stream, each regime contributes 50 construction examples.
In the imbalanced construction stream, each domain down-samples its designated minority regime to 5 examples and keeps 30 examples for each dominant regime.

\subsection{Conflict Pair Construction}

Conflict pairs are the core of RegimeBench.
For a base input $x$, we construct two or more instances with similar surface content but different active regimes:
\begin{equation}
    (x,c,r_a,y^{r_a})
    \quad \text{and} \quad
    (x,c',r_b,y^{r_b}).
\end{equation}
The output $y^{r_a}$ is correct under $r_a$ but becomes a cross-regime negative under $r_b$, and vice versa.
This yields the negative set $\mathcal{Y}_i^{-}$ used by the CRI metric.
In moderation, conflict pairs are created by changing the policy context while holding the post nearly fixed.
In robot planning, they are created by changing the operational objective while holding the environment fixed.
In cultural norm adaptation, they are created by changing the social regime while holding the communicative intent fixed.

\subsection{Minority Regime Splits}

To test minority-regime query preservation, we vary the frequency of one or more regimes during memory construction.
In v0, the minority regimes are \textit{minor safety} for moderation, \textit{energy-saving} for robot planning, and \textit{hierarchy-sensitive} for cultural norms.
The minority regime appears only 5 times in imbalanced construction streams while dominant regimes appear 30 times each; minority regimes are then oversampled in the query phase.
A robust memory system should preserve these memories despite their low frequency.
We evaluate both retained minority-memory count and performance on minority-regime queries; the main tables report the performance-based MRP-Acc value.

\subsection{Evolving Regime Splits}

In evolving-regime splits, the validity of some labels changes between construction and query.
In v0, crisis-event moderation becomes stricter, efficiency-first robot planning inherits a hard safety constraint, and informal-peer cultural responses require a softened alternative under the updated regime.
Because RAM-minimal does not implement deprecation or split/merge, evolving-regime results are diagnostic rather than the central evidence for the method.

\subsection{Output Evaluation}

For the reported v0 pilot, moderation, robot planning, and cultural norm adaptation all use exact-match evaluation over deterministic output labels or response templates.
A future LLM judge protocol for free-form generation is not used to compute the reported numbers.

\section{Metric Details}
\label{app:metric_details}

\paragraph{Regime-conditioned accuracy.}
RCA measures whether the final output matches the target behavior under the active regime.
In the reported v0 pilot, all domains use exact-match labels or deterministic response templates; regime-aware judges are reserved for future free-form generation.

\paragraph{Cross-regime interference.}
CRI uses the explicit negative set $\mathcal{Y}^{-}$.
An output counts as interference when it matches or semantically follows a behavior that is valid under another regime but invalid under the current one.

\paragraph{Homogenization rate.}
HR measures generic outputs that avoid committing to regime-specific behavior.
For instance, ``be respectful'' is homogenized when the target requires a direct refusal under one regime and deferential negotiation under another.
In v0, homogenization is detected by matching a domain-specific generic output label.

\paragraph{Minority-regime query preservation.}
The main tables report MRP-Acc: downstream accuracy on minority-regime queries.
The released JSON files keep the shorter key \texttt{MRP}; \texttt{metric\_aliases.json} maps it to the paper's MRP-Acc terminology.
The artifacts also log memory-state retention counts in \texttt{memory\_states.jsonl}.
A memory is considered retained if it remains accessible, non-deprecated, and correctly assigned or compatible with the minority regime.

\paragraph{Invalid memory activation.}
IMA is computed over the effective retrieved context.
It measures the fraction of activated memories whose assigned regime is incompatible with the current regime and which are not invariant.
IMA is especially useful for diagnosing failures hidden by final-answer accuracy.

\section{Reproducibility, Ethics, and LLM Usage}
\label{app:reproducibility_ethics_llm}

\paragraph{Reproducibility.}
The reported pilot uses deterministic templates, TF-IDF retrieval, exact-match task labels, fixed bootstrap seed 7, and no external services.
The artifact records the generated dataset, predictions, retrieved memories, memory states, metrics, diagnostics, ablations, capacity sensitivity results, memory-cost summaries, metric aliases, and failure cases under \texttt{results/regimebench\_v0/}.
The experiment scripts are under \texttt{code/regimebench\_v0/}.
The released \texttt{data\_summary.json} records 600 balanced construction prototypes, 2,085 logged split-specific construction events, 720 queries, and 2,805 total rows.
All table values are traceable to JSON artifacts; Monolithic Memory uses \texttt{N/A} for IMA in the paper because monolithic summaries do not expose individual regime assignments.
For the tagged-capacity diagnostic, RAM-minimal is first built normally and then pruned by removing the oldest memories from the largest stores until its realized total memory count matches Regime-Tagged Retrieval for the same domain and split.

\paragraph{Ethics.}
RegimeBench v0 uses synthetic moderation, robot-planning, and cultural-norm templates rather than human-subject data.
The benchmark is intended to diagnose memory validity failures, not to define real moderation, safety, or cultural policy.
Future real-world versions should include domain review and calibrated evaluation protocols.

\paragraph{LLM usage.}
Large language model assistance was used for drafting, editing, code-writing support, and consistency checks.
No LLM outputs are used as labels, predictions, judge scores, metric values, or hand-filled trends in the reported pilot.
The authors are responsible for the paper's claims, experiments, artifact contents, and final wording.


\end{document}
